\section{Detector Positioning}

The determination of the position of the far detector was determined
using the information provided by the GPS receivers installed in the
timing system masters (TDU masters) at Ash River.  These three receivers
were located in the far detector computing center in relay rack 3 and
connected to a set of three antennas that are mounted on the roof of
the Ash River building.

The position of the near detector on the surface (NDOS) was determined
through the use of the master TDU mounted in the first electronics
rack at the NOvA surface building upper catwalk and connected to the
antenna mounted on the roof of the building.  Additionally the
position of the MINOS surface building was determined in a similar
manner using the TDUs mounted in the electronics racks in the MINOS
surface building and connected to two separate antennas mounted on a
mast on the roof of the building.

\subsection{GPS Antenna Position}

The Antennas on the Ash River building are mounted on a single mast
with four mount points for Antennas.  Each Antenna is mounted on a
mount point at a distance XXX from the central axis of the mast.  The
antennas are then connected to the recievers through an approximately
100 foot cable run that is common to all the antennas and receivers.

At the MINOS service building a similar antenna installations is used
with a single mast with two mount points and common cable runs down to
the GPS receivers.  For the NOvA surface building a single mast and
antenna were mounted and only a single cable runs down from the
antenna to the receiver.

\subsection{TDU NMEA Data Stream}

To extract the calculated GPS position, the receiver's NMEA data
stream output was connected to a serial console using a USB to serial
adapter (USB/serial chip on TDU).  The console was brought up using
the GNU ``screen'' utility using a connection speed of 38400 baud.

When this was started the MNEA data stream was continously sent to the serial
console in a fashion where it repleats its NMEA block 1 once every 5
seconds.  When extracted the blocks read similar to:

\begin{verbatim}
$GPGLL,4822.716417,N,09249.877638,W,210341.002,A,A*41
$GPRMC,210341.002,A,4822.716417,N,09249.877638,W,0.000,0.00,070313,,,A*40
$GPGGA,210341.002,4822.716417,N,09249.877638,W,1,12,0.71,382.823,M,-32.462,M,,*5D
$GPGSA,A,3,31,11,14,25,01,20,23,30,22,16,12,32,1.20,0.71,0.96*0C
$GPGSV,3,1,12,31,88,142,48,11,02,241,38,14,35,092,45,25,31,053,44*75
$GPGSV,3,2,12,01,20,256,44,20,40,300,45,23,11,285,42,30,27,171,44*74
$GPGSV,3,3,12,22,07,152,42,16,02,182,37,12,07,029,40,32,64,275,47*7C
$POLYT,210341.002,070313,421421.0022814,1730,421437.0022814,-281415,333.984,25,1127648768,00017,00017,*4C
$POLYP,210341.002,4822.716417,N,09249.877638,W,350.361,G3,001.9,005.1,0.000,0.00,0.000,,0.71,0.96,1.20,1.30,0.52,12,,0*23
$POLYS,12,31,U,142,88,48,255,11,U,241,02,38,255,14,U,092,35,45,255,25,U,053,31,44,255,01,U,256,20,44,255,20,U,300,40,45,255,23,U,285,11,42,255,30,U,171,27,44,255,22,U,152,07,42,255,16,U,182,02,37,201,12,U,029,07,40,255,32,U,275,64,47,255*73
$POLYI,JN,18,EXT,HPOS,VPOS,INT,TSYNC,*19 
\end{verbatim}

The \$GPGGA string is then used to determine the position of the fix.
In the case of the data above the \$GPGGA string is decode to the
position of the Ash River antenna array as shown in Table~\ref{table:gps_decode}.

\begin{table}[thb]
\begin{center}
\begin{tabular}{| c |r | r|}
\hline
Name                     & Data                      & Description \\
\hline
Sentence Identifier & \$GPGGA              & Global Positioning System Fix Data \\
Time                      &  210341.002         & 21:03:41.002 Z \\
Latitude                 & 4822.716417,N     & 48d 22.716417' N \\
Longitude              & 09249.877638,W  & 92d 49.877638' W \\
Fix Quality             & 1                           & 0=Invalid,
1=GPS fix, 2 = DGPS fix \\
Number of Satellites & 12                      & 12 satellites in view
and locked \\
Horizontal Dilution of Precision (HDOP) & 0.71 & final std error is
0.71 times input error \\
Altitude                   & 382.823,M           & 382.823 meters
above mean sea level \\
Height of geoid above WGS84 ellipsoid & -32.462,M & -32.462 meters \\
Checksum               & \*5D & Checksum for sentence \\
\hline
\end{tabular}
\caption[GPGGA Sentence Decoding]{\label{table:gps_decode} Decoding of the \$GPGGA sentence for
the Ash River GPS antenna array.}
\end{center}
\end{table}

The \$POLYP string can be similarly decoded to provide an alternative
position fix based on a user datum ellipsoid and the \$POLYT string
provides the GPS time information along with the errors in the
measurement and associated drifts.

\section{Flight Path Determination}

The flight path between the near detector and the far detector was
computing in two different manners using different assumptions for the
shape of the Eath.  These calculations are done both for self
consistancy and to determine the potential systemmatic error that is
introduced by the model of the Earth being used.

\subsection{Spherical Earth}

In the simplest approximation the Earth can be taken as a sphere with
a mean equitoral radius $R_E = 6,378,137$ meters.  In this representation
the longitude ($\theta$), latitude ($\phi$) and altitude (A) of a point can be converted to a
cartisan coordoinates in an Earth centered reference frame
representation.

The transformation $\{\theta, \phi, A\}$ to $\{p_x,p_y,p_z\}$ is given by:

\begin{align}
p_x &= ( R_E + A ) \sin(\theta)\cos(\phi) \\ 
p_y &= (R_E + A) \cos(\theta)\cos(\phi) \\
 p_z &= (R_E + A) \sin(\phi)
\end{align} 

In this case the longitude and latitude are taken from the GPS NMEA
data stream and the altitude that is used is the altitude above mean
sea level.

In the case of the near detector to far detector the parameters sets that were used
were:

\begin{table}
\begin{center}
\begin{tabular}{| c | c | c | c| c| }
\hline
Location & Latitude & Longitude & Altitude (m) & Offset Depth to Detector (m)\\
\hline
Ash River Ant-1 & 48 d 22.715780'' & 92 d 49.876930 & 348         & 12+7.8 \\  
Ash River Ant-2 & 48 d 22.715780'' & 92 d 49.876930 & 348.638  & 12+7.8 \\  
Ash River Ant-3 & 48 d 22.716416'' & 92 d 49.877639 & 350.346  & 12+7.8 \\  
\hline
Fnal Minos         &  41 d 50 m 23 s   & 88 d 16 m 13 s  & 228
& 105 + 15 \\
\hline
\end{tabular}
\caption[GPS Detector Locations]{\label{table:gps_det_loc} GPS
  coordinates of NOvA detectors used for flight path calculations.}
\end{center}
\end{table}

Because the values returned by the GPS receivers are the positions of
the antenna systems, offsets are applied to these reported positions
to account for the difference between the location of the antenna and
the location of the detector.  In each case the offset is computed to
correspond to the front face of the detector (face closest to the
beam) so that the flight time delays that are computed are for the
earliest arrival of particles between sites (i.e. delays are computing
from front face crossing to front face crossing and from that
corrected for position).  The depth offsets are computed to the
mid-point height of the detector (7.8~m for far det) and include the
approximate height of the detector over the top edge (12~m for far
det).  For the NuMI cavern the depth of the cavern is taken to be
105~m and the mount point of the antenna on the Minos service building
is taken as 15~m above the surface.

\begin{table}
\begin{center}
\begin{tabular}{| c | c | c | c| c| }
\hline
Location & $\Delta x_{offset}$ (m) & $Delta y_{offset}$ (m) & $Delta
z_{offset}$ (m)& Notes\\
\hline
Ash River Ant Base & TBD & TBD & TBD  & mount axis \\
Ash River Ant-1     & TBD & TBD & TBD  & south ant.\\  
Ash River Ant-2     & TBD & TBD & TBD  & north ant.\\  
Ash River Ant-3     & TBD & TBD & TBD  & west ant.\\  
\hline
Fnal NDOS         &  2 & 3   & -3 &  105 + 15 \\
\hline
\end{tabular}
\caption[Antenna to Detector
Offsets]{\label{table:antenna_det_offsets} Postions offsets used to
  correct for antenna verse detector front face positions.}
\end{center}
\end{table}

To comput these offsets the plane tangent to an earth centered sphere,
at the GPS measured antenna position is taken, and the vector between
antenna position and the center of the front face of the detector is
calculated.  The projects of the offset vector into the earth centered
fixed reference frame are extracted and the earth centered fixed
coordinates of the antenna are shifted accordingly.

The standard Euclidian metric is computed between the ECF coordinates
of the the near and far detectors to determine the straight line
flight path between the detectors.

\begin{align}
\Delta L &= \sqrt{\sum_{i} ( (p_i - \delta p_i) - (p_i^\prime - \Delta
  p_i^{prime})^2} \\
             &= 810,319~\text{m}
\end{align}

The neutrinos are assumes to have a speed equal to the speed of light,
c.  The resulting flight time is then:

\begin{align}
\Delta t &= \Delta L / c\\
             &= 2703.3~\mu\text{t}
\end{align}

\subsection{WGS84 Earth Model}

Due to the lattitudes and distances between the near and far detector
sites, a more accurate determination of the flight path can found by
using a deformed Earth model.  The standard description for the
deformation can be taken from the WGS84 model which establishes
elliptical major and minor axis for the earth and a set of
transformations from the ellipsoid into an Earth Centered Fixed (ECF)
representation.

The parameters used for the WGS84 model, as used by the both the NOvA
and Minos experiments are given by:

\begin{align}
R_a &= 6378137 ~\text{m}\\
R_b &= 6356752.31424518 ~\text{m}\\
\epsilon_a &= \sqrt{\frac{R^2_a - R^2_b}{R_a}}\\
\epsilon_b &= \sqrt{\frac{R^2_a - R^2_b}{R_a}}
\end{align}

With a prescription for transforming from the GPS reported ($\theta$,
$\phi$, A) coordinates into the ECF coordinates ($p_x$,$p_y$,$p_z$).

\begin{align}
R_{curv}(\theta) &= \frac{a}{\sqrt{1-\epsilon^2_a \sin \theta^2}} \\
p_x &= (R_{curv}(\theta) + A) \cos \theta \cos \phi \\
p_y &= (R_{curv}(\theta) + A) \cos \theta \sin \phi \\
p_z &= R^2_b/R^2_a (R_{curv}(\theta) + A) \sin \theta
\end{align}

In the case of the path length between the FNAL and Ash River timing
stations the path length is taken by comparing the distance between
the ECF coordinates of far detector and the Minos service building
position (offset to the beam depth and detector cavern).
\begin{align}
\Delta L &= \sqrt{\sum_{i} ( (p_i - \delta p_i) - (p_i^\prime - \Delta
  p_i^{prime})^2} \\
             &= 809,606~\text{m}
\end{align}

The time of flight is then computed assuming the propagation of the
beam at the speed of light

\begin{align}
\Delta t &= \Delta L / c \\
             &= 2700.56~\mu \text{s}
\end{align}

This time offsite is the delay between when the beam crosses the near
detector cavern and when the beam should cross the front face of the
far detector.